% Part: first-order-logic
% Chapter: completeness
% Section: henkin-expansion

\documentclass[../../../include/open-logic-section]{subfiles}

\begin{document}

\olfileid{fol}{com}{hen}

\olsection{Henkin-uitbreiding}

\begin{explain}
Een deel van de moeilijkheid bij het bewijzen van de volledigheidsstelling
is dat het model dat we uit een volledige consistente verzameling~$\Gamma$
construeren, alle gekwantificeerde !!{formula}s in~$\Gamma$ waar moet
maken. Om dit te garanderen gebruiken we een kunstgreep van Leon Henkin.
Deze kunstgreep komt er in wezen op neer dat we de taal uitbreiden met
oneindig veel !!{constant}s en voor elke !!{formula} met één vrije
!!{variable} $!A(x)$ een formule toevoegen van de vorm
\iftag{prvEx}
      {$\lexists[x][!A(x)] \lif !A(c)$}
      {$\lnot\lforall[x][!A(x)] \lif \lnot !A(c)$},
waarbij $c$ een van de nieuwe !!{constant}s is. Wanneer we de
!!{structure} construeren die~$\Gamma$ vervult, garandeert dit dat elke
\iftag{prvEx}
{ware existentiële zin een getuige heeft}
{onware universele zin een tegenvoorbeeld heeft}
onder de nieuwe constanten.
\end{explain}

\begin{prop}
\ollabel{prop:lang-exp}
Als $\Gamma$ consistent is in $\Lang L$ en $\Lang L'$ uit $\Lang L$
wordt verkregen door !!a{denumerable} verzameling nieuwe
!!{constant}s $\Obj d_0$, $\Obj d_1$, \dots toe te voegen, dan is
$\Gamma$ consistent in~$\Lang L'$.
\end{prop}

\begin{defn}[Verzadigde verzameling]
Een verzameling $\Gamma$ van !!{formula}s van een taal $\Lang {L}$ is
\emph{verzadigd} dan en slechts dan als er voor elke
!!{formula}~$!A(x) \in \Frm[L]$ met één vrije !!{variable}~$x$ een
!!{constant}~$c \in \Lang{L}$ bestaat waarvoor
\iftag{prvEx}
      {$\lexists[x][!A(x)] \lif !A(c) \in \Gamma$}
      {$\lnot\lforall[x][!A(x)] \lif \lnot !A(c) \in \Gamma$}.
\end{defn}

De volgende definitie wordt gebruikt in het bewijs van de volgende
stelling.

\begin{defn}
\ollabel{defn:henkin-exp}
Laat $\Lang L'$ zijn zoals in \olref{prop:lang-exp}. Kies een opsomming
$!A_0(x_0)$, $!A_1(x_1)$, \dots van alle !!{formula}s~$!A_i(x_i)$
van~$\Lang L'$ waarin één variabele ($x_i$) vrij voorkomt. We definiëren
de !!{sentence}s~$!D_n$ door inductie naar~$n$.

Laat $c_0$ de eerste !!{constant} onder de $\Obj d_i$ zijn die we aan
~$\Lang{L}$ hebben toegevoegd en die niet in~$!A_0(x_0)$ voorkomt.
Veronderstel dat $!D_0$, \dots,~$!D_{n-1}$ reeds zijn gedefinieerd. Laat
$c_n$ dan de eerste van de nieuwe !!{constant}s~$\Obj d_i$ zijn die noch
in $!D_0$, \dots,~$!D_{n-1}$, noch in~$!A_n(x_n)$ voorkomt.

Laat nu $!D_{n}$ de volgende !!{formula} zijn: \iftag{prvEx}
{$\lexists[x_{n}][!A_{n}(x_{n})] \lif
  !A_{n}(c_{n})$}{$\lnot\lforall[x_{n}][!A_{n}(x_{n})] \lif \lnot
  !A_{n}(c_{n})$}.
\end{defn}

\begin{lem}
\ollabel{lem:henkin}
Elke consistente verzameling~$\Gamma$ kan worden uitgebreid tot een
verzadigde consistente verzameling~$\Gamma'$.
\end{lem}

\begin{proof}
Neem een consistente verzameling zinnen~$\Gamma$ in een taal~$\Lang{L}$.
Breid de taal uit tot~$\Lang{L'}$ door !!a{denumerable} verzameling
nieuwe !!{constant}s toe te voegen. Volgens \olref{prop:lang-exp} blijft
$\Gamma$ consistent in de rijkere taal. Laat verder $!D_i$ zijn zoals
in \olref{defn:henkin-exp}. Definieer
\begin{align*}
\Gamma_0 & = \Gamma \\
\Gamma_{n+1} & = \Gamma_n \cup \{!D_n \}
\end{align*}
dat wil zeggen, $\Gamma_{n+1} = \Gamma \cup \{ !D_0, \dots, !D_n \}$,
en stel $\Gamma' = \bigcup_{n} \Gamma_n$. Dan is $\Gamma'$ kennelijk
verzadigd.

Als $\Gamma'$ inconsistent was, dan zou voor een zekere $n$ ook
$\Gamma_n$ inconsistent zijn (Opgave: leg uit waarom). Om aan te tonen
dat $\Gamma'$ consistent is, volstaat het dus door inductie naar~$n$ te
bewijzen dat elke verzameling~$\Gamma_n$ consistent is.

De inductiebasis is eenvoudigweg de bewering dat $\Gamma_0 = \Gamma$
consistent is; dit is de hypothese van de stelling. Veronderstel voor de
inductiestap dat $\Gamma_{n}$ consistent is, maar $\Gamma_{n+1} =
\Gamma_n \cup \{!D_n\}$ inconsistent is. Bedenk dat $!D_n$ gelijk is aan
\iftag{prvEx}
{$\lexists[x_{n}][!A_{n}(x_n)] \lif !A_{n}(c_{n})$}
{$\lnot\lforall[x_{n}][!A_{n}(x_n)] \lif \lnot !A_{n}(c_{n})$},
waarbij $!A_n(x_n)$ !!a{formula} van $\Lang{L'}$ is waarin alleen de
variabele~$x_n$ vrij voorkomt. Uit de keuze van~$c_n$ (zie
\olref{defn:henkin-exp}) volgt dat $c_n$ noch in~$!A_n(x_n)$, noch
in~$\Gamma_n$ voorkomt.

Als $\Gamma_n \cup \{!D_n\}$ inconsistent is, dan $\Gamma_n
\Proves \lnot !D_n$, en dus gelden beide volgende beweringen:
\iftag{prvEx}{
\[
\Gamma_n \Proves \lexists[x_n][!A_n(x_n)]
\qquad
\Gamma_n \Proves \lnot !A_n(c_n)
\]}{
\[
\Gamma_n \Proves \lnot\lforall[x_n][!A_n(x_n)]
\qquad
\Gamma_n \Proves !A_n(c_n)
\]}
Omdat $c_n$ niet voorkomt in
$\Gamma_n$ of in~$!A_n(x_n)$, is
\tagrefs{prfAX/{fol:axd:qpr:thm:strong-generalization},prfSC/{fol:seq:qpr:thm:strong-generalization},prfND/{fol:ntd:qpr:thm:strong-generalization},prfTab/{fol:tab:qpr:thm:strong-generalization}}
van toepassing. Uit \iftag{prvEx}
{$\Gamma_n \Proves \lnot !A_n(c_n)$}
{$\Gamma_n \Proves !A_n(c_n)$}
verkrijgen we
\iftag{prvEx}
{$\Gamma_n \Proves \lforall[x_n][\lnot !A_n(x_n)]$}
{$\Gamma_n \Proves \lforall[x_n][!A_n(x_n)]$}.
Daarmee gelden zowel
\iftag{prvEx}
{$\Gamma_n \Proves \lexists[x_n][!A_n(x_n)]$ als
$\Gamma_n \Proves \lforall[x_n][\lnot !A_n(x_n)]$}
{$\Gamma_n \Proves \lnot\lforall[x_n][!A_n(x_n)]$ als
$\Gamma_n \Proves \lforall[x_n][!A_n(x_n)]$},
zodat $\Gamma_n$ zelf inconsistent is.
\iftag{prvEx}
{(Merk op dat
\iftag{prvAll}
{$\lforall[x_n][\lnot !A_n(x_n)] \Proves
  \lnot\lexists[x_n][!A_n(x_n)]$}
{$\lforall[x_n][\lnot !A_n]$ is gedefinieerd als
  $\lnot\lexists[x_n][\lnot\lnot !A_n(x_n)]$ en
  $\lnot\lexists[x_n][\lnot\lnot !A_n(x_n)] \Proves
  \lnot\lexists[x_n][!A_n(x_n)]$}.)}{}
Dit is een tegenspraak, want $\Gamma_n$ was verondersteld consistent te
zijn. Bijgevolg is $\Gamma_n \cup \{ !D_n\}$ consistent.
\end{proof}

\begin{explain}
We zullen nu aantonen dat \emph{volledige}, consistente verzamelingen die
verzadigd zijn de volgende eigenschap hebben: \iftag{prvAll}{zo'n
  verzameling bevat een universeel gekwantificeerde !!{sentence} dan en
  slechts dan als zij al haar instanties bevat}{}\iftag{defAll,defEx}{}{
  en }\iftag{prvAll}{zij bevat een existentieel gekwantificeerde
  !!{sentence} dan en slechts dan als zij ten minste één instantie
  bevat}{}. We zullen dit gebruiken om aan te tonen dat de !!{structure}
die we uit een volledige, consistente, verzadigde verzameling genereren,
al haar gekwantificeerde zinnen waar maakt.
\end{explain}

\begin{prop}\ollabel{prop:saturated-instances}
Veronderstel dat $\Gamma$ volledig, consistent en verzadigd is.
\begin{tagenumerate}{prvEx,prvAll}
\tagitem{prvEx}{$\lexists[x][!A(x)] \in \Gamma$ dan en slechts dan als
  $!A(t) \in \Gamma$ voor ten minste één gesloten term~$t$.}{}
\tagitem{prvAll}{$\lforall[x][!A(x)] \in \Gamma$ dan en slechts dan als
  $!A(t) \in \Gamma$ voor alle gesloten termen~$t$.}{}
\end{tagenumerate}
\end{prop}

\begin{proof}
  \begin{tagenumerate}{prvEx,prvAll}
  \tagitem{prvEx}{%
    \iftag{probEx}{Opgave.}{Veronderstel eerst dat
      $\lexists[x][!A(x)] \in \Gamma$. Omdat $\Gamma$ verzadigd is,
      geldt $(\lexists[x][!A(x)] \lif !A(c)) \in \Gamma$ voor een
      zekere !!{constant}~$c$. Uit
      \tagrefs{prfAX/{fol:axd:ppr:prop:provability-lif},prfSC/{fol:seq:ppr:prop:provability-lif},prfND/{fol:ntd:ppr:prop:provability-lif},prfTab/{fol:tab:ppr:prop:provability-lif}}, onderdeel~(1),
      en \olref[ccs]{prop:ccs}\olref[ccs]{prop:ccs-prov-in} volgt dat
      $!A(c) \in \Gamma$.

    Voor de andere richting is verzadiging niet nodig. Veronderstel dat
    $!A(t) \in \Gamma$. Dan volgt uit
      \tagrefs{prfAX/{fol:axd:qpr:prop:provability-quantifiers},prfSC/{fol:seq:qpr:prop:provability-quantifiers},prfND/{fol:ntd:qpr:prop:provability-quantifiers},prfTab/{fol:tab:qpr:prop:provability-quantifiers}}, onderdeel~(1),
    dat $\Gamma \Proves \lexists[x][!A(x)]$. Volgens
    \olref[ccs]{prop:ccs}\olref[ccs]{prop:ccs-prov-in} geldt dan
    $\lexists[x][!A(x)] \in \Gamma$.}}{}
  \tagitem{prvAll}{%
    \iftag{probAll}{Opgave.}{Veronderstel dat $!A(t) \in \Gamma$ voor
      alle gesloten termen~$t$. Neem ter tegenspraak aan dat
      $\lforall[x][!A(x)] \notin \Gamma$. Omdat $\Gamma$ volledig is,
      geldt $\lnot\lforall[x][!A(x)] \in \Gamma$. Uit verzadiging volgt
      dat \iftag{prvEx}{$(\lexists[x][\lnot !A(x)] \lif \lnot !A(c)) \in
        \Gamma$}{$(\lnot\lforall[x][!A(x)] \lif \lnot !A(c)) \in
        \Gamma$} voor een zekere !!{constant}~$c$. Omdat $c$ een gesloten
      term is, geldt volgens de aanname $!A(c) \in \Gamma$. Maar dan zou
      $\Gamma$ inconsistent zijn. (Opgave: geef de !!{derivation} waaruit
      blijkt dat
      \[
      \lnot \lforall[x][!A(x)],
      \iftag{prvEx}{\lexists[x][\lnot !A(x)] \lif \lnot
        !A(c)}{\lnot\lforall[x][!A(x)] \lif \lnot
        !A(c)}, !A(c)
      \]
      inconsistent is.)

      Voor de omgekeerde richting hebben we verzadiging niet nodig.
      Veronderstel dat $\lforall[x][!A(x)] \in \Gamma$. Dan volgt uit
      \tagrefs{prfAX/{fol:axd:qpr:prop:provability-quantifiers},prfSC/{fol:seq:qpr:prop:provability-quantifiers},prfND/{fol:ntd:qpr:prop:provability-quantifiers},prfTab/{fol:tab:qpr:prop:provability-quantifiers}},
      onderdeel~(2), dat $\Gamma \Proves !A(t)$. Volgens
      \olref[ccs]{prop:ccs} geldt $!A(t) \in \Gamma$.}}{}
  \end{tagenumerate}
\end{proof}

\end{document}
